% Original course diagram. All text, mathematics, and connectors remain vector.
\begin{figure}[htbp]
\centering
\begin{tikzpicture}[
  font=\small,
  flow/.style={-{Latex[length=2mm]},thick,draw=navy},
  target/.style={-{Latex[length=2mm]},draw=navy,dashed},
  update/.style={-{Latex[length=2mm]},thick,draw=coral!85!black},
  box/.style={draw=navy,fill=navy!5,rounded corners=2pt,
    align=center,minimum height=0.88cm,inner sep=4pt},
  train/.style={box,draw=teal!80!black,fill=teal!8},
  check/.style={box,draw=navy,fill=codebg}
]
\node[box,text width=1.5cm] (corpus) at (1,0) {Text\\corpus};
\node[box,text width=1.8cm] (split) at (3.25,0) {Split text\\before windows};
\node[train,text width=2.1cm] (trainwin) at (6,1.3)
  {Training\\windows $(x,y)$};
\node[check,text width=2.1cm] (valwin) at (6,-1.3)
  {Validation\\windows $(x,y)$};
\node[train,text width=1.9cm] (trainmodel) at (9,1.3)
  {GPT $\theta$\\logits $z$};
\node[check,text width=1.9cm,minimum height=1.08cm] (valmodel) at (9,-1.3)
  {GPT $\theta$\\no gradients\\logits $z$};
\node[train,text width=1.4cm] (trainloss) at (11.7,1.3)
  {$\mathrm{CE}(z,y)$\\training loss};
\node[check,text width=1.4cm] (valloss) at (11.7,-1.3)
  {$\mathrm{CE}(z,y)$\\held-out loss};
\node[train,text width=1.55cm,draw=coral!85!black,fill=coral!8]
  (optim) at (14.05,1.3) {Backward\\AdamW step};
\node[align=center,text width=1.65cm,text=navy] (report) at (14.05,-1.3)
  {Report and\\compare};

\draw[flow] (corpus) -- (split);
\draw[flow] (split.east) -- ++(0.35,0) |- (trainwin.west);
\draw[flow] (split.east) -- ++(0.35,0) |- (valwin.west);
\draw[flow] (trainwin) -- node[above] {$x$} (trainmodel);
\draw[flow] (valwin) -- node[above] {$x$} (valmodel);
\draw[flow] (trainmodel) -- (trainloss);
\draw[flow] (valmodel) -- (valloss);
\draw[flow] (trainloss) -- (optim);
\draw[flow] (valloss) -- (report);

% Shifted target IDs go directly to the loss, not into the decoder input.
\draw[target] (trainwin.south) -- (6,0.2)
  -- node[below] {targets $y$} (11.7,0.2) -- (trainloss.south);
\draw[target] (valwin.south) -- (6,-2.4)
  -- node[below] {targets $y$} (11.7,-2.4) -- (valloss.south);
\draw[update] (optim.north) -- (14.05,2.6)
  -- node[above] {update $\theta$; repeat with the next batch} (9,2.6)
  -- (trainmodel.north);
\end{tikzpicture}
\caption{The pretraining data and optimization paths. Split the corpus before
forming overlapping windows, tokenize each split with the same fixed tokenizer,
and shift each window by one position to obtain input IDs $x$ and target IDs $y$.
Cross-entropy (CE) compares the logits with the targets. Only the training branch
computes gradients and updates parameters. Validation uses the same current
parameters in evaluation mode, with dropout disabled and no gradient recording.
If the tokenizer is learned from the course corpus, fit it only on the training
split.}
\label{fig:pretraining-flow}
\end{figure}
